%Basic constructions like size and language :
\documentclass[12pt]{report}
\usepackage[utf8]{inputenc}
\usepackage{CJKutf8} %Chinese, support of begin cjk end cjk
\usepackage[T1]{fontenc} %French
\usepackage[letterpaper, margin=1.27cm]{geometry}
\setlength{\emergencystretch}{3em} %the very wide text block needs slack to avoid overfull lines
\usepackage{subfiles} 
\usepackage{times} %Times New Roman font

%References:
\usepackage[round]{natbib} %author-year citations, loaded before hyperref

%Make title, table of contents
\usepackage{hyperref} %necessay for maketitle to work
\hypersetup{colorlinks=true, allcolors=blue} %all links and citations in blue
\usepackage{datetime2} %for the dtmnow command

\title{Casual Document for 
Quantitative Mixing Time Results for Single Site Gibbs in Stochastic Block Model}
\author{Wen, Zehai; Cesare, Michael}
\date{\DTMnow}

%Page style, header and footer:
\usepackage{fancyhdr}

\pagestyle{fancy}
\setlength{\headheight}{14.5pt}
\addtolength{\topmargin}{-2.5pt}
\fancyhf{}
\lhead{
Wen, Zehai; Cesare, Michael}
\rhead{\today}

\rfoot{}
\cfoot{\thepage}

%Graphics and figures:
\usepackage{graphicx} 
\usepackage{float} %both graphicx and float packages are necessary for begin figure
\usepackage{caption} %necessary for any kind of caption
\usepackage{subcaption} %necessary for subfigure

%Tabular:
\usepackage{multirow}

%Coding:
\usepackage{listings}
\usepackage{color}

\lstset{literate=%
  {é}{{\'e}}{1}%
  {è}{{\`e}}{1}%
  {à}{{\`a}}{1}%
  {ç}{{\c{c}}}{1}%
  {œ}{{\oe}}{1}%
  {ù}{{\`u}}{1}%
  {É}{{\'E}}{1}%
  {È}{{\`E}}{1}%
  {À}{{\`A}}{1}%
  {Ç}{{\c{C}}}{1}%
  {Œ}{{\OE}}{1}%
  {Ê}{{\^E}}{1}%
  {ê}{{\^e}}{1}%
  {î}{{\^i}}{1}%
  {ô}{{\^o}}{1}%
  {û}{{\^u}}{1}%
  {ä}{{\"{a}}}1
  {ë}{{\"{e}}}1
  {ï}{{\"{i}}}1
  {ö}{{\"{o}}}1
  {ü}{{\"{u}}}1
  {û}{{\^{u}}}1
  {â}{{\^{a}}}1
  {Â}{{\^{A}}}1
  {Î}{{\^{I}}}1
} 

\definecolor{dkgreen}{rgb}{0,0.6,0}
\definecolor{gray}{rgb}{0.5,0.5,0.5}
\definecolor{mauve}{rgb}{0.58,0,0.82}

\lstset{frame=tb,
  language=Java,
  aboveskip=3mm,
  belowskip=3mm,
  showstringspaces=false,
  columns=flexible,
  basicstyle={\small\ttfamily},
  numbers=none,
  numberstyle=\tiny\color{gray},
  keywordstyle=\color{blue},
  commentstyle=\color{dkgreen},
  stringstyle=\color{mauve},
  breaklines=true,
  breakatwhitespace=true,
  tabsize=3
}

%Graph Theory:
\usepackage{tikz}
\usetikzlibrary{arrows.meta}

%Commands :
\newcommand\tab[1][1cm]{\hspace*{#1}}
\renewcommand{\bibname}{Références}

%%%%%%%%%%%%%%%%%%%%%%%%%%%%Mathematical packages%%%%%%%%%%%%%%%%%%%%%%%%%%%%%%%%%%%%%%%%%%

%Basic utilities:
\usepackage{amssymb}
\usepackage{amsmath}
\usepackage{amsfonts}
\usepackage{physics} %for vector symbol
\usepackage{bm} %for bm

%Commands :
\DeclareMathOperator*{\argmax}{arg\,max}
\DeclareMathOperator*{\argmin}{arg\,min}

\newcommand{\bracket}[1]{\left ( #1 \right )}
\newcommand{\curly}[1]{\left \{ #1 \right \} }
\newcommand{\squared}[1]{\left [ #1 \right ]}

\usepackage{amsthm}  %Theorem styles
\theoremstyle{plain}
\newtheorem{theorem}{Theorem}

\theoremstyle{definition}
\newtheorem{definition}[theorem]{Definition} 
\newtheorem{definitions}[theorem]{Definitions} 
\newtheorem{example}[theorem]{Example}
\newtheorem{examples}[theorem]{Examples}
\newtheorem{remark}[theorem]{Remark}
\newtheorem{remarks}[theorem]{Remarks}
\newtheorem{lemma}[theorem]{Lemma}
\newtheorem{corollary}[theorem]{Corollary}
\newtheorem{proposition}[theorem]{Proposition}

%%%%%%%%%%%%%%%%%%%%%%%%%%%%%%%%%%%%%%%%%%%%%%%%%%%%%%%%%%%%%%%%%%%%%%%%%%%%%%%%

\begin{document}
\begin{CJK*}{UTF8}{gbsn}
\maketitle

\setcounter{tocdepth}{0}
\tableofcontents

\chapter{The Stochastic Block Model and the Single Site Gibbs Method}

\section{Basic Setup}\label{sec:basic_setup}

The procedure of analyzing graphical data 
using Bayesian stochastic block model with single 
site Gibbs procedure seems to have been first 
mentioned in \cite{Kemp_et_al_2004}
and slightly later in \cite{Xu_et_al_2006}.
In this section, we
describe the simplest setup of both the model and the single site Gibbs method
by following a tutorial article \cite{Schmidt_Morup_2013}.

For the moment, 
let us consider the simplest case where
the graphical data received by the statistician is just
one simple graph of $N \geqslant 3$ vertices
labeled as $1, 2, \ldots, N$,
represented by its adjacency matrix $X \in \{0,1\}^{N \times N}$
where $X_{i,j} = 1$ if and only if there is an edge between vertex $i$ and vertex $j$.
In particular, we have $X_{i,i} = 0$ for all $i$ and
$X_{i,j} = X_{j,i}$ for all $i,j$.
If we want to, 
we can consider later the case which our data contains 
multiple graphs of the same size generated independently identically distributed 
from a distribution in the same model, but we shall not do so in this section.

We assume that $X$ is generated from a distribution 
in the (standard Bayesian) stochastic block model (SBM)
with unknown number of blocks. 
To describe the model, 
let us write $\Pi$
for the collection of all partitions of the set $\{1, 2, \ldots, N\}$
and represent
any element $z \in \Pi$ using the \textit{restricted growth string} notation,
meaning that we see $z$ 
as a vector $z = (z_1, z_2, \ldots, z_N)$
where $z_i$ denotes the block to which vertex $i$ belongs,
such that the classes are ordered as they first appear: $z_1 = 1$
and, for every $i \in \{2, 3, \ldots, N\}$, we have

\begin{equation*}
  z_i \in \curly{1, 2, \ldots, 1 + \max_{j < i} z_j}.
\end{equation*}
For example, if $N = 5$ and the partition is $z = \{\{1, 3\}, \{2\}, \{4, 5\}\}$,
then we write $z = (1, 2, 1, 3, 3)$.
As derived in \cite{Schmidt_Morup_2013},
(the standard Bayesian)
SBM is a parameteric model for the matrix $X$ whose only
parameter is $z$, ranged over $\Pi$,
with density function is given by 

\begin{equation}\label{eq:SBM_density}
  p(X \mid z) =
  \prod_{1 \leqslant k \leqslant l \leqslant K} 
  \frac{B(m_{k,l}+a, \bar{m}_{k,l}+b)}{B(a,b)},
\end{equation}
where $a > 0$ and $b > 0$ are hyperparameters fixed by the 
statistician, $B(\cdot, \cdot)$ is the beta function,
$K$ is the number of non-empty blocks as specified by $z$,
$m_{k,l}$ is the number of edges between block $k$ and block $l$,
and $\bar{m}_{k,l}$ is the number of non-edges between block $k$ and block $l$.

(Actually, just as in \cite{Schmidt_Morup_2013}, 
we have abused the notation a bit in \eqref{eq:SBM_density}.
Given that $X$ is symmetric with diagonal entries equal to $0$,
we are only interested in the vector $(X_{i,j})_{1 \leqslant i < j \leqslant N}$
and we have written $p(X \mid z)$ 
in place of $p((X_{i,j})_{1 \leqslant i < j \leqslant N} \mid z)$.)

In order to do Bayesian statistics, 
we have to specify a prior distribution for $z$.
As pointed out in \cite{Schmidt_Morup_2013},
we will use the Chinese restaurant process (CRP) prior for $z$.
With a hyperparameter $A > 0$ fixed by the statistician,
$\text{CRP}(A)$ is a probability measure on $\Pi$ defined as 
follows: $z_1 \equiv 1$ and, for every $i \in \{2, 3, \ldots, N\}$,
given $z_1, z_2, \ldots, z_{i-1}$ who have already used a total of 
$K^{(i-1)}$ blocks, 
$z_i$ will join an existing block $k \in \{1, 2, \ldots, K^{(i-1)}\}$ with probability

\begin{equation*}
  \frac{n_k^{(i-1)}}{A + i - 1},
\end{equation*}
where $n_k^{(i-1)}$ is the number of vertices among $1, 2, \ldots, i-1$
who have been assigned to block $k$,
and will create a new block $K^{(i-1)} + 1$ with probability

\begin{equation*}
  \frac{A}{A + i - 1}.
\end{equation*}
$\text{CRP}(A)$ does have a clean joint density (with respect to the counting measure on $\Pi$)

\begin{equation}\label{eq:CRP_density}
 p(z) = \frac{A^{K} \Gamma(A)}{\Gamma(A + N)} \prod_{k=1}^{K} \Gamma(n_k),
\end{equation}
where $K$ is the number of non-empty blocks as specified by $z$,
$n_k$ is the number of vertices assigned to block $k$ according to $z$,
and $\Gamma(\cdot)$ is the gamma function.

By Bayes's theorem, according to \eqref{eq:SBM_density}
and \eqref{eq:CRP_density},
the posterior distribution of $z$ given $X$ is

\begin{equation*}
  p(z \mid X) \propto p(X \mid z) p(z)
  = \squared{
  \prod_{1 \leqslant k \leqslant l \leqslant K}
  \frac{B(m_{k,l}+a, \bar{m}_{k,l}+b)}{B(a,b)}}
  \squared{\frac{A^{K} \Gamma(A)}{\Gamma(A + N)} \prod_{k=1}^{K} \Gamma(n_k)}.
\end{equation*}
We now move to carry out the single site Gibbs method.
For every $i \in \{1, 2, \ldots, N\}$,
let $z^{(-i)}$ denote the vector $(z_1, z_2, \ldots, z_{i-1}, z_{i+1}, \ldots, z_N)$
obtained from $z$ by removing the $i$-th entry.
We need to compute, for every $i \in \{1, 2, \ldots, N\}$,
what's proportional to

\begin{equation*}
  p(z_i \mid z^{(-i)}, X).
\end{equation*}
Therefore, for each $k \in \{1, 2, \ldots, K^{(-i)}\}$,
where $K^{(-i)}$ is the number of non-empty blocks as specified by $z^{(-i)}$,
we must calculate $p(z_i = k \mid z^{(-i)}, X)$.

\section{Gibbs Draft}

In this section we explain how we understand the conditional distribution
$p(z_i \mid z^{(-i)}, X)$ used by the single site Gibbs method,
given that $z$ ranges over the set of partitions $\Pi$
rather than over a product space such as $\mathbb{R}^N$.
Throughout the section,
we fix a vertex $i \in \{1, 2, \ldots, N\}$
and we regard every $z \in \Pi$ as a set of blocks,
each block being a non-empty subset of $\{1, 2, \ldots, N\}$.
The restricted growth string notation is only used to write down
a partition after it has been updated;
it plays no role in the update itself.

\subsection*{The partition of the other vertices and where vertex $i$ goes}

Let $\Pi^{(-i)}$ denote the collection of all partitions of the set
$\{1, 2, \ldots, N\} \setminus \{i\}$.
Every $z \in \Pi$ induces a partition of the other $N - 1$ vertices,
namely its restriction

\begin{equation*}
  T_i(z) = \curly{C \setminus \{i\} : C \in z} \setminus \curly{\varnothing}
  \in \Pi^{(-i)}.
\end{equation*}
For $\rho \in \Pi^{(-i)}$ with $K_\rho$ blocks $B_1, B_2, \ldots, B_{K_\rho}$,
which we name by their vertex sets rather than by integer labels,
let

\begin{equation*}
  \mathcal{C}_\rho = \curly{B_1, B_2, \ldots, B_{K_\rho}, \mathrm{new}}
\end{equation*}
be the set of the $K_\rho + 1$ options for vertex $i$:
join one of the existing blocks of $\rho$, or start a new block.
For $c \in \mathcal{C}_\rho$, we write $\rho \oplus_i c \in \Pi$
for the partition obtained from $\rho$ by putting $i$ into $c$, that is,

\begin{equation*}
  \rho \oplus_i B_k = \bracket{\rho \setminus \{B_k\}} \cup \curly{B_k \cup \{i\}}
  \quad \text{for } k \in \{1, 2, \ldots, K_\rho\},
  \qquad
  \rho \oplus_i \mathrm{new} = \rho \cup \curly{\{i\}}.
\end{equation*}

\begin{lemma}\label{lem:gibbs_bijection}
The map $(\rho, c) \mapsto \rho \oplus_i c$ is a bijection from
$\curly{(\rho, c) : \rho \in \Pi^{(-i)},\ c \in \mathcal{C}_\rho}$
onto $\Pi$.
Its inverse sends $z \in \Pi$ to $(T_i(z), c_i(z))$,
where, writing $C$ for the block of $z$ containing $i$,
we set $c_i(z) = C \setminus \{i\}$ if $C \neq \{i\}$
and $c_i(z) = \mathrm{new}$ if $C = \{i\}$.
\end{lemma}

\begin{proof}
For every $\rho \in \Pi^{(-i)}$ and $c \in \mathcal{C}_\rho$,
removing $i$ from $\rho \oplus_i c$ gives back $\rho$,
and the block of $\rho \oplus_i c$ containing $i$ is $c \cup \{i\}$
if $c \neq \mathrm{new}$ and $\{i\}$ if $c = \mathrm{new}$,
so $T_i(\rho \oplus_i c) = \rho$ and $c_i(\rho \oplus_i c) = c$.
Conversely, for every $z \in \Pi$,
putting $i$ back into $c_i(z)$ recovers $z$,
so $T_i(z) \oplus_i c_i(z) = z$.
\end{proof}

In the coordinates given by Lemma~\ref{lem:gibbs_bijection},
$\rho = T_i(z)$ plays the role of $z^{(-i)}$
and $c = c_i(z)$ plays the role of $z_i$.
Accordingly, for $\rho \in \Pi^{(-i)}$ and $c \in \mathcal{C}_\rho$, we define

\begin{equation}\label{eq:gibbs_conditional}
  p(z_i = c \mid z^{(-i)} = \rho, X)
  = \frac{p(\rho \oplus_i c \mid X)}
  {\sum_{c' \in \mathcal{C}_\rho} p(\rho \oplus_i c' \mid X)}
  \propto p(X \mid \rho \oplus_i c)\, p(\rho \oplus_i c).
\end{equation}
This is a probability distribution over the $K_\rho + 1$ options in $\mathcal{C}_\rho$,
not over integer labels: no block label appears anywhere.
If $i$ forms a block on its own in the current state $z$,
then $T_i(z)$ has one block fewer than $z$
and the current state corresponds to the option $c = \mathrm{new}$;
this case needs no special treatment.

\begin{remark}
The single site Gibbs update for vertex $i$ is the Markov kernel on $\Pi$ given by

\begin{equation*}
  P_i(z, z') = \frac{p(z' \mid X)}{\sum_{z'' \in \Pi : T_i(z'') = T_i(z)} p(z'' \mid X)}
  \, \mathbf{1}\curly{T_i(z') = T_i(z)},
  \qquad z, z' \in \Pi,
\end{equation*}
that is, we keep $\rho = T_i(z)$ fixed and redraw $c$
according to \eqref{eq:gibbs_conditional}.
If $T_i(z) = T_i(z')$, then both $p(z \mid X) P_i(z, z')$ and
$p(z' \mid X) P_i(z', z)$ are equal to $p(z \mid X) p(z' \mid X)$
divided by the same normalizing sum,
and otherwise both vanish.
Hence $P_i$ is reversible with respect to $p(\cdot \mid X)$.
No coordinates, linear structure or topology on $\Pi$ are needed for this.
\end{remark}

\subsection*{Vertex $i$ as the last vertex to arrive}

For a permutation $\sigma$ of $\{1, 2, \ldots, N\}$,
write $\sigma(z) = \curly{\sigma(C) : C \in z}$
and let $X^\sigma$ be the adjacency matrix given by
$X^\sigma_{\sigma(j), \sigma(l)} = X_{j,l}$.
The CRP density \eqref{eq:CRP_density} depends on $z$ only through
$K$ and the block sizes, so $p(\sigma(z)) = p(z)$.
Likewise, the density \eqref{eq:SBM_density} depends on $(X, z)$ only through
the numbers of edges and non-edges between blocks,
which do not change when the vertices of $X$ and $z$ are permuted together,
so $p(X^\sigma \mid \sigma(z)) = p(X \mid z)$.
Hence the posterior does not care about the order of the vertices,
and in \eqref{eq:gibbs_conditional} we may pretend that $i$ is the $N$-th vertex
to arrive, even when $i = 1$.
The computations below confirm this picture:
the conditional distribution \eqref{eq:gibbs_conditional} factors as

\begin{equation*}
  p(z_i = c \mid z^{(-i)} = \rho, X)
  \propto
  \squared{\text{CRP seating probability of the last vertex}}
  \times
  \squared{\text{how well the edges of $i$ fit $c$}}.
\end{equation*}

\paragraph{Prior part.}
For every block $B$ of $\rho$, write $n_B = |B|$.
Since $\Gamma(n_B + 1) = n_B \Gamma(n_B)$ and $\Gamma(1) = 1$,
\eqref{eq:CRP_density} gives

\begin{equation*}
  p(\rho \oplus_i B) = \frac{A^{K_\rho} \Gamma(A)}{\Gamma(A + N)}
  \prod_{B' \in \rho} \Gamma(n_{B'}) \cdot n_B,
  \qquad
  p(\rho \oplus_i \mathrm{new}) = \frac{A^{K_\rho} \Gamma(A)}{\Gamma(A + N)}
  \prod_{B' \in \rho} \Gamma(n_{B'}) \cdot A,
\end{equation*}
where the second identity uses that $\rho \oplus_i \mathrm{new}$
has $K_\rho + 1$ blocks.
Hence

\begin{equation*}
  p(\rho \oplus_i c) \propto
  \begin{cases}
    n_B & \text{if } c = B \in \rho, \\
    A & \text{if } c = \mathrm{new}.
  \end{cases}
\end{equation*}
Since $\sum_{B \in \rho} n_B = N - 1$,
the normalized weights are $n_B / (A + N - 1)$ and $A / (A + N - 1)$,
which are exactly the CRP probabilities for the $N$-th vertex
to join block $B$ or to create a new block.

\paragraph{Likelihood part.}
Write $X^{(-i)} = (X_{j,l})_{1 \leqslant j < l \leqslant N,\ j, l \neq i}$
for the graph with vertex $i$ removed
and $x_i = (X_{i,j})_{j \neq i}$ for the edges of vertex $i$.
For blocks $B, B'$ of $\rho$ (with $B = B'$ allowed),
let $m_{B,B'}$ and $\bar{m}_{B,B'}$ be the numbers of edges and non-edges
between $B$ and $B'$ in $X^{(-i)}$,
where each unordered pair of distinct vertices is counted once
(so that, when $B = B'$, these are the pairs of vertices inside $B$).
For every block $B'$ of $\rho$, let

\begin{equation*}
  r_{B'} = \sum_{j \in B'} X_{i,j}
\end{equation*}
be the number of edges between $i$ and $B'$,
so that there are $n_{B'} - r_{B'}$ non-edges between $i$ and $B'$.
Finally, let

\begin{equation*}
  p(X^{(-i)} \mid \rho) = \prod_{\{B, B'\}}
  \frac{B(m_{B,B'} + a, \bar{m}_{B,B'} + b)}{B(a,b)},
\end{equation*}
where the product ranges over all unordered pairs $\{B, B'\}$ of blocks of $\rho$
with $B = B'$ allowed;
this is the right-hand side of \eqref{eq:SBM_density}
applied to the graph $X^{(-i)}$ and the partition $\rho$.

\begin{proposition}\label{prop:gibbs_likelihood}
For every block $B$ of $\rho$,

\begin{equation*}
  p(X \mid \rho \oplus_i B) = p(X^{(-i)} \mid \rho)
  \prod_{B' \in \rho}
  \frac{B(m_{B,B'} + r_{B'} + a, \bar{m}_{B,B'} + n_{B'} - r_{B'} + b)}
  {B(m_{B,B'} + a, \bar{m}_{B,B'} + b)},
\end{equation*}
and

\begin{equation*}
  p(X \mid \rho \oplus_i \mathrm{new}) = p(X^{(-i)} \mid \rho)
  \prod_{B' \in \rho}
  \frac{B(r_{B'} + a, n_{B'} - r_{B'} + b)}{B(a,b)}.
\end{equation*}
\end{proposition}

\begin{proof}
The blocks of $\rho \oplus_i B$ are $B \cup \{i\}$ and the blocks
$B' \neq B$ of $\rho$.
Pairs of vertices not involving $i$ are counted exactly as in $\rho$.
For every $B' \in \rho$, the $n_{B'}$ pairs $\{i, j\}$ with $j \in B'$
fall into the block pair $\{B \cup \{i\}, B'\}$ if $B' \neq B$,
and into the block pair $\{B \cup \{i\}, B \cup \{i\}\}$ if $B' = B$;
in both cases they add $r_{B'}$ edges and $n_{B'} - r_{B'}$ non-edges
to the counts $m_{B,B'}$ and $\bar{m}_{B,B'}$.
So, in \eqref{eq:SBM_density}, only the factors of the block pairs
containing $B \cup \{i\}$ change,
there is exactly one such factor for each $B' \in \rho$,
and dividing each new factor by the corresponding factor of $p(X^{(-i)} \mid \rho)$
gives the first identity.
The blocks of $\rho \oplus_i \mathrm{new}$ are $\{i\}$ and the blocks of $\rho$.
The block pairs of $\rho$ keep their counts,
the block pair $\{\{i\}, B'\}$ receives $r_{B'}$ edges and $n_{B'} - r_{B'}$ non-edges
for every $B' \in \rho$,
and the block pair $\{\{i\}, \{i\}\}$ contains no pair of vertices,
so its factor is $B(a,b)/B(a,b) = 1$.
This gives the second identity.
\end{proof}

The factor $p(X^{(-i)} \mid \rho)$ is the same for every $c \in \mathcal{C}_\rho$
and therefore drops out of \eqref{eq:gibbs_conditional}.
Since the marginal distribution of $X^{(-i)}$ under $\rho \oplus_i c$
is the SBM with partition $\rho$
(it only involves pairs of vertices not containing $i$),
the remaining product is the conditional probability
$p(x_i \mid X^{(-i)}, \rho \oplus_i c)$ of the edges of $i$ given the rest of the graph.
Each of its factors is a Beta--Bernoulli posterior predictive probability:
in the hierarchical formulation of \cite{Schmidt_Morup_2013},
the link probability $\theta_{B,B'}$ between blocks $B$ and $B'$
has prior $\text{Beta}(a,b)$,
its posterior given $X^{(-i)}$ is $\text{Beta}(m_{B,B'} + a, \bar{m}_{B,B'} + b)$,
and

\begin{equation*}
  \frac{B(m_{B,B'} + r_{B'} + a, \bar{m}_{B,B'} + n_{B'} - r_{B'} + b)}
  {B(m_{B,B'} + a, \bar{m}_{B,B'} + b)}
\end{equation*}
is the probability, under this posterior, of the particular configuration of
$r_{B'}$ edges and $n_{B'} - r_{B'}$ non-edges between $i$ and $B'$.
For a new block there is no previous observation,
so the posterior is simply the prior $\text{Beta}(a,b)$.
These are exactly the quantities computed (on the logarithmic scale)
in the code of \cite{Schmidt_Morup_2013}.

\paragraph{Conclusion.}
Combining the prior part with Proposition~\ref{prop:gibbs_likelihood},
for every $\rho \in \Pi^{(-i)}$ we obtain

\begin{equation}\label{eq:gibbs_conditional_explicit}
  p(z_i = c \mid z^{(-i)} = \rho, X) \propto
  \begin{cases}
    \displaystyle n_B \prod_{B' \in \rho}
    \frac{B(m_{B,B'} + r_{B'} + a, \bar{m}_{B,B'} + n_{B'} - r_{B'} + b)}
    {B(m_{B,B'} + a, \bar{m}_{B,B'} + b)}
    & \text{if } c = B \in \rho, \\[4ex]
    \displaystyle A \prod_{B' \in \rho}
    \frac{B(r_{B'} + a, n_{B'} - r_{B'} + b)}{B(a,b)}
    & \text{if } c = \mathrm{new},
  \end{cases}
\end{equation}
where the proportionality constant is fixed by summing over
the $K_\rho + 1$ options in $\mathcal{C}_\rho$.
In words, $p(z_i \mid z^{(-i)}, X)$ is the probability that vertex $i$
joins each block of $\rho$, or starts a new one,
given the partition $\rho$ of the other vertices:
it is the CRP seating weight ($n_B$ or $A$)
times the posterior predictive likelihood of the edges of $i$ under that choice,
normalized over the $K_\rho + 1$ options.
This description never uses $\mathbb{R}^N$ or block labels.
After drawing $c$, we write the new state $\rho \oplus_i c$
in restricted growth string notation.

\bibliographystyle{plainnat}
\def\bibname{References} % Nom obligatoire de la section des références.
                                              % On utilise \'e car le é cause des problèmes
                                              % dans la table des matière
%% ENGLISH
 %\def\bibname{References}
\bibliography{reference.bib}     % La base de données contenant des entrées bibliographiques.
                                    % Seules celles référencées dans le texte seront ajoutées
                                    % \`a la bibliographie.

\end{CJK*}
\end{document}
